A Figura~\ref{fig:modelo} apresenta o modelo do eProfile.

\begin{figure}
  \centering\includegraphics[width=0.5\textwidth]{imagens/modelo.png}
  \caption{Modelo do eProfile}
  \label{fig:modelo}
\end{figure}

A figura tradicional do usuário na personalização é extendida aqui pelo conceito de \emph{entidade}. Uma entidade é semelhante a um usuário, porém em um conceito mais amplo, não representando apenas pessoas mas também objetos que possam criar uma trilha e ter um perfil, como um veículo ou um dispositivo móvel. Uma entidade pode também ser composta por um grupo de entidades como, por exemplo, uma frota de caminhões onde cada caminhão tem seu perfil único, e onde a frota tem um perfil também.

Os \emph{clientes} são os aplicativos que são utilizados pelas entidades para se conectar aos gerencidadores de trilhas e perfil. Neste modelo, o cliente se conecta a um gerenciador de trilhas genérico, que é responsável pelo armazenamento da trilha da entidade. O perfil da entidade será inferido a partir das trilhas registradas no gerenciador.

As trilhas das entidades são buscadas do gerenciador de trilhas pelo eProfile. Como o eProfile pode se conectar a diferentes gerenciadores de trilha, onde em cada um deles a trilha tem um formato específico, é necessário que seja realizada uma conversão entre o formato original e o formato utilizado pelo eProfile, que será apresentado mais adiante. Este trabalho é realizado por um agente entitulado \emph{conversor}, que é capaz de realizar a conversão para os formados de trilha suportados pelo eProfile.

Desta forma, o eProfile não tem a necessidade de conhecer o cliente -- mas precisa conhecer o formato utilizado pelo gerenciador de trilhas.

A trilha do usuário é então armazenada dentro do eProfile em uma \emph{ontologia de trilha básica}. Esta ontologia possui as classes básicas para a construção de uma trilha, sendo que a mesma pode ser extendida de acordo com a especificação do gerenciador de trilhas que está sendo utilizado. A forma como ocorre esta extensão é detalhada na Seção~\ref{sec:ontologias}.
